%% ============================================================
%%  Handout — Aula 13: MCTS e a família AlphaZero
%%  Prof. Marcelo Keese Albertini — UFU
%%  Compilar: latexmk -xelatex handout-aula13.tex
%% ============================================================
\documentclass[a4paper, 11pt]{article}
%% .sty do projeto na raiz do repositório (../)
\makeatletter
\def\input@path{{./}{../}}
\makeatother


\usepackage{handoutAprendizadoRL}
\usepackage{babel}
\babelprovide[import, main]{brazilian}
\usepackage{multicol}

\setaula{Aula 13}
\setprof{Prof. Marcelo Keese Albertini}

% árvore reutilizável para as fases do MCTS
\newcommand{\mctsarvore}[1]{%
\begin{tikzpicture}[baseline=(r.north), scale=0.9, transform shape]
  \tikzset{no/.style={circle, draw=rlAzul, fill=rlAzul!10, thick,
                      minimum size=4.4mm, inner sep=0pt}}
  \node[no] (r) at (0,0) {};
  \node[no] (a) at (-0.9,-0.8) {};
  \node[no] (b) at (0.9,-0.8) {};
  \node[no] (c) at (-1.4,-1.6) {};
  \node[no] (d) at (-0.4,-1.6) {};
  \draw[trans] (r)--(a);  \draw[trans] (r)--(b);
  \draw[trans] (a)--(c);  \draw[trans] (a)--(d);
  #1
\end{tikzpicture}}

\newcommand{\Mmod}{\mathcal{M}}
\newcommand{\Nv}[1]{N(#1)}
\newcommand{\Qh}{\widehat{Q}}
\newcommand{\Vh}{\widehat{V}}
\newcommand{\ret}{G}
\newcommand{\raiz}{s_{\mathrm{raiz}}}
\newcommand{\folha}{s_{\mathrm{folha}}}

\begin{document}

\capa{Busca em árvore de Monte Carlo e a família AlphaZero}
     {Busca baseada em simulação \textbullet\ UCT \textbullet\ AlphaGo, AlphaZero e MuZero}
     {Prof.\ Marcelo Keese Albertini \textbullet\ Faculdade de Computação}
     {Adaptado e expandido de \emph{CS234: Reinforcement Learning}, Emma Brunskill, Stanford University --- \emph{Lecture 13}}

\begin{ideiabox}
Todo o resto do curso constrói objetos \emph{globais}: uma tabela $\Vstar$, uma
rede $\pol_\theta$. Esta aula faz a pergunta oposta --- se eu tenho um simulador
e alguns segundos sobrando, o que consigo comprar de qualidade de decisão
\emph{aqui, neste estado}? E, mais importante: como converter esse esforço
efêmero em aprendizado permanente. A resposta a essa segunda pergunta é o
algoritmo AlphaZero.
\end{ideiabox}

\section{Planejamento em tempo de decisão}

Vale separar dois usos de um modelo:

\begin{itemize}
  \item \textbf{Planejamento em segundo plano} (\emph{background planning}).
        O modelo é usado para melhorar uma política ou função valor
        \emph{global}, que depois é consultada em tempo de execução. Iteração
        de valor e Dyna são disso.
  \item \textbf{Planejamento em tempo de decisão} (\emph{decision-time
        planning}). O modelo é usado \emph{no momento de agir}, a partir do
        estado corrente, e o resultado da computação é descartado logo depois.
        MCTS é disso.
\end{itemize}

O segundo é atraente quando $\abs{\gS}$ é grande demais para qualquer objeto
global ser bom em todo lugar, mas o agente visita, em uma partida, uma fração
minúscula dos estados. Concentrar o esforço nessa fração é uma economia
gigantesca.

\begin{defbox}{Modelo generativo}
Uma caixa-preta $\Mmod$ que, dado $(s,a)$, devolve amostras
$s'\sim \tran{\cdot}{s,a}$ e $r \sim \rec(\cdot\mid s,a)$. Não exige acesso às
probabilidades, nem enumeração de sucessores. Em jogos de tabuleiro, o modelo
\emph{são} as regras: perfeito e barato.
\end{defbox}

\begin{alertabox}
Todo o método depende dessa hipótese. Em domínios com modelo aprendido, o erro
por passo se compõe ao longo do horizonte simulado: um simulador com $1\%$ de
erro por passo é praticamente ruído após cinquenta passos. É por isso que MCTS
brilhou em jogos e continua difícil em robótica de contato ou em saúde.
\end{alertabox}

\section{Busca de Monte Carlo simples}

\begin{algbox}{Busca de Monte Carlo simples}
Dados $\Mmod$, uma política de simulação $\pol$ e um orçamento $K$ por ação:
\begin{enumerate}
  \item para cada $a\in\gA$, simule $K$ episódios a partir do estado real $s_t$
        começando com a ação $a$ e seguindo $\pol$ depois;
  \item estime $\displaystyle \Qh(s_t,a) = \frac1K\sum_{k=1}^K \ret^k_t$;
  \item execute $a_t = \argmax_a \Qh(s_t,a)$.
\end{enumerate}
\end{algbox}

\begin{teobox}{Proposição --- um passo de melhoria}
Seja $\pol'$ a política que joga $\argmax_a \Qpi(s_t,a)$ em $s_t$ e segue $\pol$
daí em diante. Então $V^{\pol'}(s_t) \ge \Vpi(s_t)$, e a busca de Monte Carlo
simples converge para $\pol'$ quando $K\to\infty$.
\end{teobox}

\begin{provabox}
Por definição, $\Qh(s_t,a) \to \Qpi(s_t,a)$ pela lei dos grandes números, pois
cada episódio é uma amostra independente do retorno sob $\pol$ após a ação $a$.
Logo o $\argmax$ empírico converge para o $\argmax$ verdadeiro (assumindo
unicidade). E
\[
  V^{\pol'}(s_t) = \max_a \Qpi(s_t,a)
  \;\ge\; \sum_a \pol(a\mid s_t)\, \Qpi(s_t,a) = \Vpi(s_t),
\]
porque o máximo domina qualquer média. \hfill$\square$
\end{provabox}

Duas consequências práticas. Primeiro, o \textbf{teto}: por melhor que seja a
amostragem, o resultado nunca supera $\pol'$ --- um único passo de melhoria
sobre a política de simulação. Segundo, o \textbf{ruído}: $\Qh$ é média de
retornos, cuja variância cresce com o horizonte. Tomar o $\argmax$ de
estimativas ruidosas produz \emph{viés de maximização}: em expectativa,
$\E[\max_a \Qh] \ge \max_a \E[\Qh]$, isto é, superestimamos sistematicamente o
valor da ação escolhida --- o mesmo fenômeno que motivou o \emph{double}
Q-learning na Aula 4.

\subsection{A árvore expectimax e por que ela não serve}

A alternativa ideal é maximizar em \emph{todos} os níveis, não só no primeiro:
\[
  Q_h(s,a) \;=\; \rec(s,a) \;+\; \gamma \sum_{s'} \tran{s'}{s,a}\,
  \max_{a'} Q_{h-1}(s',a'), \qquad Q_0 \equiv 0,
\]
avaliado por busca dirigida para frente a partir de $s_t$ --- resolver apenas o
\emph{sub-MDP} que começa agora, e não o MDP inteiro. O problema é o tamanho:
uma árvore completa de profundidade $H$ tem $(\abs{\gS}\abs{\gA})^H$ nós. Para o
\emph{Go}, com ramificação $b\approx 250$ e partidas de $d\approx 150$ jogadas,
isso está fora de cogitação, e a poda alfa-beta --- que exige minimax
determinístico e uma boa função de avaliação escrita à mão --- apenas divide o
expoente por dois.

\begin{ideiabox}
MCTS ataca as duas maldições ao mesmo tempo: \textbf{amostra} em vez de somar
sobre $s'$ (largura) e \textbf{trunca} a subárvore abaixo de uma folha,
substituindo-a por uma estimativa de valor (profundidade).
\end{ideiabox}

\section{Busca em árvore de Monte Carlo}

\subsection{As quatro fases}

\begin{center}
\begin{tabular}{@{}cccc@{}}
  \mctsarvore{\draw[trans destac] (r)--(a); \draw[trans destac] (a)--(d);}
  &
  \mctsarvore{\draw[trans destac] (r)--(a); \draw[trans destac] (a)--(d);
              \node[circle, draw=rlLaranja, fill=rlLaranja!25, thick,
                    minimum size=4.4mm, inner sep=0pt] (e) at (-0.4,-2.4) {};
              \draw[trans destac] (d)--(e);}
  &
  \mctsarvore{\node[circle, draw=rlLaranja, fill=rlLaranja!25, thick,
                    minimum size=4.4mm, inner sep=0pt] (e) at (-0.4,-2.4) {};
              \draw[trans] (d)--(e);
              \draw[rlVerde!70!black, thick, decorate,
                    decoration={snake, amplitude=0.5mm, segment length=2.2mm},
                    -{Stealth[length=2mm]}] (e) -- ++(0,-0.8);
              \node[font=\tiny, color=rlVerde!60!black] at (0.45,-3.2) {$\ret$};}
  &
  \mctsarvore{\node[circle, draw=rlLaranja, fill=rlLaranja!25, thick,
                    minimum size=4.4mm, inner sep=0pt] (e) at (-0.4,-2.4) {};
              \draw[{Stealth[length=2mm]}-, very thick, rlVerde!70!black] (d)--(e);
              \draw[{Stealth[length=2mm]}-, very thick, rlVerde!70!black] (a)--(d);
              \draw[{Stealth[length=2mm]}-, very thick, rlVerde!70!black] (r)--(a);}
  \\[0.5em]
  {\scriptsize\bfseries\color{rlAzul} 1. Seleção} &
  {\scriptsize\bfseries\color{rlAzul} 2. Expansão} &
  {\scriptsize\bfseries\color{rlAzul} 3. Simulação} &
  {\scriptsize\bfseries\color{rlAzul} 4. Retropropagação}
\end{tabular}
\end{center}

\begin{enumerate}
  \item \textbf{Seleção.} A partir da raiz, desça pela árvore escolhendo ações
        com a \emph{política de árvore} (UCT ou PUCT), que depende das
        estatísticas acumuladas. Pare ao chegar a um nó com alguma ação nunca
        tentada, ou a um estado terminal.
  \item \textbf{Expansão.} Acrescente à árvore o nó alcançado.
  \item \textbf{Simulação.} Estime o valor da nova folha. Ou por
        \emph{rollout} --- jogar até o fim com uma política fixa e barata --- ou
        por uma rede de valor $v_\theta(\folha)$.
  \item \textbf{Retropropagação.} Percorra o caminho de volta até a raiz
        atualizando $\Nv{s,a}$, $W(s,a)$ e $\Qh(s,a)$.
\end{enumerate}

A separação entre a \emph{política de árvore} (adaptativa, cara, informada) e a
\emph{política de rollout} (fixa, barata, ignorante) é o que dá ao algoritmo o
seu perfil: ele investe onde já sabe alguma coisa e economiza onde não sabe.

\subsection{Estatísticas mantidas}

\begin{center}
\begin{tabular}{@{}lll@{}}
  \toprule
  \textbf{Estatística} & \textbf{Significado} & \textbf{Atualização} \\
  \midrule
  $\Nv{i,a}$ & vezes que $a$ foi escolhida no nó $i$ & $+1$ por visita \\
  $W(i,a)$   & soma dos retornos observados          & $+\,\ret$ \\
  $\Qh(i,a)$ & valor médio $W(i,a)/\Nv{i,a}$         & recalculado \\
  $\Nv{i}$   & $\sum_a \Nv{i,a}$                     & $+1$ \\
  \bottomrule
\end{tabular}
\end{center}

A memória cresce com o número de nós \emph{visitados}, não com $\abs{\gS}$.
É esse detalhe --- e não a esperteza da regra de seleção --- que torna o método
aplicável a espaços de estados astronômicos.

\subsection{UCT}

Dentro da árvore, cada nó é um problema de bandit: os braços são as ações, o
ganho é o retorno da simulação que passa por ali. Aplicar UCB1 em cada nó dá o
\textbf{UCT}.

\begin{teobox}{Regra de seleção UCT (Kocsis e Szepesvári, 2006)}
\[
  a \;=\; \argmax_{a\in\gA(i)}\;
  \underbrace{\Qh(i,a)}_{\text{aproveitamento}}
  \;+\;
  \underbrace{c\sqrt{\frac{\ln \Nv{i}}{\Nv{i,a}}}}_{\text{exploração}},
  \qquad
  \Qh(i,a)=\frac{1}{\Nv{i,a}}\sum_{k=1}^{\Nv{i,a}} \ret^k(i,a).
\]
Ações com $\Nv{i,a}=0$ recebem bônus infinito e são visitadas primeiro.
\end{teobox}

Três observações que costumam ser fonte de erro:

\begin{itemize}
  \item o bônus decai como $\Nv{i,a}^{-1/2}$ e cresce apenas como
        $\sqrt{\ln \Nv{i}}$ --- dobrar o orçamento total quase não muda o
        equilíbrio, o que faz a busca \emph{concentrar} esforço;
  \item $c$ só tem significado se os retornos estiverem em uma escala conhecida.
        Normalize $\ret$ para $[0,1]$ (ou $[-1,1]$ em jogos de soma zero);
  \item como as estatísticas mudam a cada simulação, a política de árvore muda
        junto: o algoritmo faz melhoria de política \emph{durante} a busca, e não
        apenas ao final.
\end{itemize}

\begin{algbox}{UCT}
{\small
\begin{tabbing}
xx\=xx\=xx\=xx\=\kill
\textbf{entrada:} $s_t$, modelo $\Mmod$, orçamento $K$, constante $c$ \\
árvore $\gets$ raiz $\raiz = s_t$ \\
\textbf{para} $k=1,\dots,K$: \\
\> $s \gets \raiz$; \; caminho $\gets [\,]$ \\
\> \textbf{enquanto} $s$ está na árvore e todas as ações de $s$ já foram tentadas: \\
\>\> $a \gets \argmax_a \Qh(s,a) + c\sqrt{\ln \Nv{s}/\Nv{s,a}}$ \\
\>\> $(s',r) \sim \Mmod(s,a)$; \; anexe $(s,a,r)$ ao caminho; \; $s \gets s'$ \\
\> \textbf{se} $s$ não é terminal: acrescente $s$ à árvore \\
\> $\ret \gets \textsc{Avalia}(s)$ \\
\> \textbf{para} $(s,a,r)$ no caminho, de trás para frente: \\
\>\> $\ret \gets r + \gamma\ret$; \;\;
     $\Nv{s,a} \mathrel{+}= 1$; \;\;
     $W(s,a) \mathrel{+}= \ret$; \;\;
     $\Qh(s,a) \gets W(s,a)/\Nv{s,a}$ \\
\textbf{devolva} $\argmax_a \Nv{\raiz,a}$
\end{tabbing}}
\end{algbox}

\begin{ideiabox}
\textbf{Por que devolver a contagem de visitas, e não $\Qh$?}
Porque $\Nv{\raiz,a}$ é uma estatística acumulada sobre toda a busca: uma ação
só recebe muitas visitas se pareceu boa \emph{de forma persistente}. Já
$\Qh(\raiz,a)$ pode ser inflado por um único retorno afortunado em um ramo
pouco visitado --- exatamente o caso em que a variância é maior.
\end{ideiabox}

\subsection{O que se pode provar --- e o que não}

\begin{teobox}{Consistência do UCT}
Sob retornos limitados e exploração suficiente, a probabilidade de o UCT
selecionar uma ação subótima na raiz tende a zero quando o número de simulações
$K\to\infty$, e a estimativa na raiz converge para o valor ótimo do sub-MDP.
\end{teobox}

\begin{alertabox}
A garantia é assintótica e o pior caso é severo: Coquelin e Munos (2007)
exibem problemas em que a convergência do UCT é \emph{duplamente exponencial na
profundidade}. A intuição é simples --- o bônus logarítmico do UCB pressupõe
que as recompensas de um braço são estacionárias, o que não vale dentro de uma
árvore, onde o valor de um ramo muda à medida que a subárvore é explorada.
Domínios com uma única sequência estreita de jogadas boas (``agulha no
palheiro'') derrotam o algoritmo.
\end{alertabox}

Na prática o método funciona por dois motivos que a teoria de pior caso não
captura: os domínios costumam ter muitos caminhos razoáveis em vez de um só, e
--- sobretudo --- um bom \emph{prior} de política reduz drasticamente a
ramificação efetiva. É esse segundo ponto que a seção seguinte explora.

\subsection{Vantagens, limitações e remédios}

\begin{multicols}{2}
{\small
\textbf{Vantagens}
\begin{itemize}
  \item busca best-first altamente seletiva;
  \item avalia dinamicamente o estado atual;
  \item amostragem quebra a maldição da dimensão;
  \item funciona com modelo caixa-preta;
  \item \emph{anytime} e paralelizável;
  \item dispensa heurística de domínio;
  \item mais computação $\Rightarrow$ melhor decisão, sem retreinar.
\end{itemize}

\columnbreak

\textbf{Limitações e remédios}
\begin{itemize}
  \item ramificação alta $\to$ \emph{prior} de política (PUCT);
  \item ações contínuas $\to$ \emph{progressive widening};
  \item \emph{rollouts} fracos $\to$ rede de valor $v_\theta$;
  \item poucos dados por nó $\to$ RAVE/AMAF, transposições;
  \item modelo imperfeito $\to$ horizonte curto, ou modelo aprendido no espaço
        de valores (MuZero);
  \item recompensa esparsa $\to$ horizonte maior ou \emph{shaping}.
\end{itemize}}
\end{multicols}

\section{De AlphaGo a MuZero}

\subsection{AlphaGo (2016)}

Três redes e uma busca: uma rede de política treinada por aprendizado
supervisionado sobre partidas humanas ($p_\sigma$), uma rede de política
refinada por auto-jogo ($p_\rho$), e uma rede de valor $v_\theta$ treinada sobre
posições de auto-jogo. A busca usava $p_\sigma$ como \emph{prior} e avaliava
folhas por uma mistura
\[
  \Vh(\folha) = (1-\lambda)\, v_\theta(\folha) + \lambda\, z_{\text{rollout}},
\]
combinando duas fontes com erros de natureza diferente. Um detalhe
contraintuitivo e instrutivo: a rede \emph{supervisionada} funcionava melhor
como \emph{prior} da busca do que a rede refinada por RL --- porque é mais
\emph{diversa}, e uma busca precisa de diversidade para não colapsar.

\subsection{AlphaGo Zero (2017): remoção sucessiva}

\begin{center}
{\small
\begin{tabular}{@{}lcc@{}}
  \toprule
  & \textbf{AlphaGo} & \textbf{AlphaGo Zero} \\
  \midrule
  Partidas humanas            & sim, 30M posições & \color{rlVerde!70!black}\bfseries não \\
  \emph{Features} de domínio  & sim, dezenas      & \color{rlVerde!70!black}\bfseries só o tabuleiro \\
  \emph{Rollouts} na busca    & sim               & \color{rlVerde!70!black}\bfseries não \\
  Redes                       & três              & \color{rlVerde!70!black}\bfseries uma, duas cabeças \\
  Arquitetura                 & convolucional     & residual \\
  \bottomrule
\end{tabular}}
\end{center}

A rede única produz $(\bm{p}, v) = f_\theta(s)$, com $\bm{p}\in\Delta(\gA)$ o
\emph{prior} de política e $v\in[-1,1]$ o valor da posição do ponto de vista do
jogador da vez.

\begin{teobox}{Regra de seleção PUCT}
\[
  a \;=\; \argmax_{a}\; Q(s,a) \;+\; c_{\text{puct}}\, P(s,a)\,
  \frac{\sqrt{\sum_b \Nv{s,b}}}{1+\Nv{s,a}},
  \qquad P(s,a) = p_a .
\]
\end{teobox}

Comparada ao UCT: o bônus é \emph{multiplicado} pelo \emph{prior}, de modo que
ações implausíveis para a rede praticamente não são exploradas; o decaimento
$1/(1+\Nv{s,a})$ é mais agressivo que $1/\sqrt{\Nv{s,a}}$; e não há logaritmo.
Com $b\approx250$, é o \emph{prior} que torna a busca viável --- ela examina a
sério meia dúzia de jogadas em vez de duzentas e cinquenta. A rede não substitui
a busca: ela a poda.

\subsection{A ideia central: a busca é um operador de melhoria de política}

Ao fim de $K$ simulações, defina a \textbf{política da busca} na raiz:
\[
  \pol_{\text{MCTS}}(a\mid s) \;=\;
  \frac{\Nv{s,a}^{1/\tau}}{\sum_b \Nv{s,b}^{1/\tau}}, \qquad \tau>0 .
\]

Empiricamente, $\pol_{\text{MCTS}}$ joga melhor que o \emph{prior} $\bm{p}$ que a
gerou --- a busca gastou computação e converteu isso em qualidade. AlphaZero
fecha o ciclo tratando $\pol_{\text{MCTS}}$ como \emph{alvo de treinamento}
para $\bm{p}$, e o resultado da partida $z$ como alvo para $v$.

\begin{center}
\begin{tikzpicture}
  \def\rr{2.1}
  \node[caixa, draw=rlLaranja, fill=rlLaranja!10, text width=2.5cm, inner sep=4pt]
    (busca) at (0,\rr) {\textbf{MCTS}\\ \color{rlCinza}guiada por $f_\theta$,
                        produz $\pol_{\text{MCTS}}$};
  \node[caixa, draw=rlAzulClaro, fill=rlAzulClaro!8, text width=2.5cm, inner sep=4pt]
    (jogo) at (\rr+1.5,0) {\textbf{Auto-jogo}\\ \color{rlCinza}$a_t\sim\pol_{\text{MCTS}}$,
                           resultado $z$};
  \node[caixa, draw=rlTeal, fill=rlTeal!8, text width=2.5cm, inner sep=4pt]
    (dados) at (0,-\rr) {\textbf{Buffer}\\ \color{rlCinza}triplas
                         $(s_t, \pol_{\text{MCTS}}, z)$};
  \node[caixa, draw=rlAzul, fill=rlAzul!8, text width=2.5cm, inner sep=4pt]
    (treino) at (-\rr-1.5,0) {\textbf{Treino de $f_\theta$}\\
                              \color{rlCinza}$\bm{p}\to\pol_{\text{MCTS}}$, $v\to z$};
  \draw[trans destac] (busca) to[bend left=22] (jogo);
  \draw[trans destac] (jogo) to[bend left=22] (dados);
  \draw[trans destac] (dados) to[bend left=22] (treino);
  \draw[trans destac] (treino) to[bend left=22] (busca);
  \node[font=\scriptsize\itshape, color=rlCinza, align=center] at (0,0)
    {nenhum dado humano\\ entra neste ciclo};
\end{tikzpicture}
\end{center}

\begin{teobox}{Perda do AlphaGo Zero / AlphaZero}
Para cada posição $s$ do auto-jogo, com alvos $\bm{\pi}=\pol_{\text{MCTS}}(\cdot\mid s)$
e $z\in\{-1,0,+1\}$:
\[
  \ell(\theta) \;=\;
  \termo{(z - v_\theta(s))^2}{regressão de valor}
  \;-\;
  \termo{\bm{\pi}^{\!\top}\log \bm{p}_\theta(s)}{destilação da busca}
  \;+\;
  \termo{c\norm{\theta}^2}{regularização}.
\]
\end{teobox}

Esta é a leitura que vale guardar: \textbf{é iteração de política}, com a busca
no papel do passo de melhoria e o auto-jogo no papel da avaliação. Anthony,
Tian e Barber (2017) batizaram o esquema de \emph{expert iteration} --- a busca
é o ``especialista lento'', a rede é o ``aprendiz rápido'' que \emph{amortiza}
o especialista, tornando a próxima busca mais barata e mais afiada. Note também
o que \emph{não} aparece na perda: nenhum gradiente de política, nenhuma razão
de importância, nenhuma região de confiança. Todo o problema difícil de RL foi
empurrado para dentro da busca, e o que sobrou é aprendizado supervisionado.

\subsection{Duas engenhocas indispensáveis}

\begin{defbox}{Ruído de Dirichlet na raiz}
$P(\raiz,a) = (1-\varepsilon)\,p_a + \varepsilon\,\eta_a$, com
$\bm{\eta}\sim\mathrm{Dir}(\alpha)$ e $\varepsilon\approx0{,}25$. Garante que
toda jogada legal tenha alguma chance de ser explorada, mesmo que a rede a
considere impossível.
\end{defbox}

\begin{defbox}{Temperatura}
Nas primeiras jogadas de cada partida usa-se $\tau=1$ (amostragem proporcional
às visitas), o que gera diversidade de aberturas; depois $\tau\to0$
(praticamente $\argmax$), o que garante jogo forte.
\end{defbox}

As duas são a versão de auto-jogo do problema de exploração das Aulas 10--12.
Sem elas, o sistema converge para um repertório estreito e passa a treinar
contra um adversário degenerado --- o auto-jogo perde a propriedade que o torna
valioso, que é gerar automaticamente um currículo de dificuldade crescente.

\subsection{AlphaZero e MuZero}

\textbf{AlphaZero (2018)} aplicou a mesma receita, sem conhecimento de domínio
além das regras, a \emph{Go}, xadrez e shogi, superando os melhores programas
especializados de cada um. Avaliava ordens de grandeza \emph{menos} posições por
jogada que os motores clássicos de xadrez --- evidência direta de que um bom
julgamento posicional vale por muitas simulações.

\textbf{MuZero (2020)} removeu a última peça de conhecimento humano: as regras.
Aprende três funções --- representação $h_\theta$ (observação $\to$ estado
latente), dinâmica $g_\theta$ ($(z,a)\to(z',\hat r)$) e predição $f_\theta$
($z\to(\bm{p},v)$) --- e roda a busca inteiramente no espaço latente. Não há
decodificador: $z$ não precisa reconstruir o estado real, apenas predizer bem
\emph{recompensa, valor e política}, que é tudo o que o planejador consome. O
modelo é aprendido \emph{para servir ao planejamento}, e não para ser fiel.

\begin{discbox}
Se o estado latente do MuZero não precisa reconstruir o mundo, em que sentido
ele ainda é um ``modelo''? O que se ganha e o que se perde ao abandonar a
exigência de fidelidade --- e como isso se relaciona com a discussão de
especificação de recompensa que veremos na próxima aula?
\end{discbox}

\section{O que generaliza, e o que não}

\begin{multicols}{2}
{\small
\textbf{\color{rlVerde!70!black}Transfere bem}
\begin{itemize}
  \item busca como operador de melhoria; rede como amortização da busca;
  \item \emph{prior} de política para podar ramificação alta;
  \item trocar computação de decisão por qualidade de decisão;
  \item aprender o modelo no espaço do que o planejador consome.
\end{itemize}

\columnbreak

\textbf{\color{rlVermelho}Não transfere de graça}
\begin{itemize}
  \item simulador perfeito, determinístico e barato;
  \item recompensa terminal não ambígua ($\pm1$);
  \item auto-jogo como fonte gratuita de currículo;
  \item episódios reinicializáveis à vontade.
\end{itemize}}
\end{multicols}

O laço ``gerar candidatos com um \emph{prior}, avaliar com um verificador,
destilar o resultado de volta no \emph{prior}'' reaparece hoje em modelos de
linguagem que buscam sobre passos de raciocínio. A analogia é boa, mas a peça
difícil migrou: em jogos, o verificador são as regras; em texto, não existe
regra que declare vitória, e construir esse verificador é o problema, não a
busca.

\begin{discbox}
MCTS otimiza muito bem aquilo que recebe como recompensa. Em que situações essa
competência é um risco em vez de uma virtude? Traga um exemplo de um domínio em
que você confiaria mais em uma política aprendida sem busca do que em uma busca
profunda sobre um modelo aproximado.
\end{discbox}

\needspace{0.32\textheight}
\section{Formulário}

\begin{multicols}{2}
{\small
\textbf{Busca de Monte Carlo simples}
\[ \Qh(s_t,a)=\tfrac1K\textstyle\sum_k \ret^k_t \;\to\; \Qpi(s_t,a) \]

\textbf{Backup expectimax}
\[ Q_h(s,a)=\rec(s,a)+\gamma\!\!\sum_{s'}\!\tran{s'}{s,a}\,V_{h-1}(s') \]
\[ V_{h-1}(s')=\max_{a'}Q_{h-1}(s',a') \]

\textbf{UCT}
\[ a=\argmax_a \Qh(i,a)+c\sqrt{\tfrac{\ln \Nv{i}}{\Nv{i,a}}} \]

\textbf{Retropropagação}
\[ \ret \gets r+\gamma\ret, \quad \Qh \gets \tfrac{W+\ret}{N+1} \]

\columnbreak

\textbf{PUCT}
\[ a=\argmax_a Q(s,a)+c_{\text{puct}}P(s,a)\tfrac{\sqrt{\sum_b \Nv{s,b}}}{1+\Nv{s,a}} \]

\textbf{Política da busca}
\[ \pol_{\text{MCTS}}(a\mid s)\propto \Nv{s,a}^{1/\tau} \]

\textbf{Perda AlphaZero}
\[ \ell=(z-v)^2-\bm{\pi}^{\!\top}\log\bm{p}+c\norm{\theta}^2 \]

\textbf{Ruído na raiz}
\[ P\gets(1-\varepsilon)\bm{p}+\varepsilon\bm{\eta},\; \bm{\eta}\sim\mathrm{Dir}(\alpha) \]}
\end{multicols}

\section{Exercícios}

\begin{exbox}{13.1 --- Aritmética do UCT}
Na raiz há três ações, com as estatísticas abaixo e $\Nv{\raiz}=100$, $c=1$.
\begin{center}
{\small\begin{tabular}{@{}lccc@{}}
\toprule
 & $a_1$ & $a_2$ & $a_3$ \\
\midrule
$\Nv{\raiz,a}$ & $80$ & $15$ & $5$ \\
$\Qh(\raiz,a)$ & $0{,}62$ & $0{,}55$ & $0{,}70$ \\
\bottomrule
\end{tabular}}
\end{center}
(a) Calcule o escore UCT de cada ação e diga qual seria simulada em seguida.
(b) Qual ação o algoritmo devolveria como jogada real, e por quê ela difere do
item (a)? (c) Quantas visitas $a_3$ precisaria acumular para que seu escore
caísse abaixo do de $a_1$, mantendo $\Qh$ fixo?
\end{exbox}

\begin{exbox}{13.2 --- Viés de maximização}
Sejam $\Qh(a_1)$ e $\Qh(a_2)$ estimativas não viesadas e independentes de dois
valores verdadeiros \emph{iguais}, $q_1=q_2=0$, cada uma $\mathcal{N}(0,\sigma^2)$.
Mostre que $\E[\max_a \Qh(a)] = \sigma/\sqrt{\pi} > 0 = \max_a \E[\Qh(a)]$.
Explique por que isso significa que a busca de Monte Carlo simples
\emph{superestima} sistematicamente o valor da ação escolhida, e por que o
efeito piora com o horizonte.
\end{exbox}

\begin{exbox}{13.3 --- Visitas contra valor}
Construa um exemplo numérico pequeno com duas ações na raiz em que
$\argmax_a \Nv{\raiz,a} \neq \argmax_a \Qh(\raiz,a)$, com a ação de maior $\Qh$
tendo apenas duas visitas. Argumente qual das duas regras você preferiria e
formule a preferência em termos de erro-padrão da média.
\end{exbox}

\begin{exbox}{13.4 --- PUCT e ramificação efetiva}
Suponha $\abs{\gA}=250$ e um \emph{prior} $\bm{p}$ que concentra massa $0{,}9$
em $5$ ações, distribuindo o resto uniformemente. (a) Com $\sum_b \Nv{s,b}=800$
e $c_{\text{puct}}=1{,}5$, compare o bônus PUCT de uma ação do grupo
concentrado (com $\Nv{s,a}=100$) e de uma ação da cauda (com $\Nv{s,a}=0$).
(b) Repita com $\bm{p}$ uniforme. (c) O que o resultado diz sobre a afirmação
``a rede poda a busca''?
\end{exbox}

\begin{exbox}{13.5 --- Alvos de treinamento}
Explique por que AlphaZero usa o resultado final da partida $z$ como alvo de
valor, e não o $\Qh(\raiz)$ produzido pela busca, que aparentemente é mais
informativo. Discuta em termos de viés e variância, e do risco de o valor
aprendido realimentar seus próprios erros.
\end{exbox}

\begin{exbox}{13.6 --- Sinal em jogos de dois jogadores}
Um aluno implementa MCTS para o jogo da velha e esquece de inverter o sinal do
retorno a cada nível na retropropagação. Descreva com precisão o comportamento
resultante do agente: contra que tipo de adversário ele parecerá competente, e
em que momento a falha se revela?
\end{exbox}

\begin{exbox}{13.7 --- Modelo imperfeito}
Um simulador aprendido erra o estado seguinte com probabilidade $\epsilon$ por
passo, de forma independente. (a) Qual a probabilidade de uma trajetória
simulada de $H$ passos estar inteiramente correta? (b) Para $\epsilon=0{,}02$,
qual o maior $H$ com probabilidade de acerto acima de $0{,}5$? (c) Que
implicação isso tem para a escolha entre um horizonte de busca longo com
\emph{rollouts} e um horizonte curto com rede de valor?
\end{exbox}

\section{Gabarito comentado (13.1 e 13.2)}

{\small
\textbf{13.1.} Com $\ln 100 \approx 4{,}605$: o bônus é
$\sqrt{4{,}605/80}=0{,}240$, $\sqrt{4{,}605/15}=0{,}554$ e
$\sqrt{4{,}605/5}=0{,}960$. Os escores ficam
$a_1: 0{,}860$, $a_2: 1{,}104$, $a_3: 1{,}660$ --- a próxima simulação vai para
$a_3$, a ação \emph{menos} visitada e de maior $\Qh$. Já a \emph{jogada real}
seria $a_1$, com $80$ visitas: a regra de seleção é otimista de propósito,
enquanto a regra de saída é conservadora de propósito. Para (c), precisamos de
$0{,}70 + \sqrt{4{,}605/N} < 0{,}62 + 0{,}240 = 0{,}860$, isto é
$\sqrt{4{,}605/N} < 0{,}160$, ou seja $N > 180$ --- mais visitas do que o
orçamento inteiro do exemplo. Ou seja: com $\Qh$ superior, $a_3$ não perde a
disputa por bônus; ela só perderia se seu $\Qh$ caísse ao ser mais explorada
--- que é exatamente o que se espera que aconteça, e a razão de o método
funcionar.

\smallskip
\textbf{13.2.} Para $X_1,X_2$ i.i.d.\ $\mathcal{N}(0,\sigma^2)$ vale
$\E[\max(X_1,X_2)] = \sigma/\sqrt{\pi}$ (caso particular de
$\E[\max] = \sigma\,\tfrac{n}{\sqrt\pi}$ para $n=2$ braços, obtido de
$\max(x,y) = \tfrac{x+y}{2}+\tfrac{|x-y|}{2}$ e
$\E|X_1-X_2| = 2\sigma/\sqrt{\pi}$). Como $\max_a \E[\Qh(a)]=0$, o excesso é
puro artefato do ruído. Na busca de Monte Carlo simples, $\sigma$ é o
erro-padrão da média de $K$ retornos, que cresce com a variância do retorno ---
e a variância do retorno cresce com o horizonte. Logo o viés otimista é maior
justamente onde a busca é mais necessária. É esse mecanismo que motiva
estimadores desacoplados (\emph{double} Q-learning) e a preferência por
contagem de visitas na saída do MCTS.}

\section{Glossário}

{\footnotesize
\begin{tabular}{@{}p{3.3cm}p{4.15cm}p{3.15cm}p{4.15cm}@{}}
\toprule
\textbf{Inglês} & \textbf{Português} & \textbf{Inglês} & \textbf{Português} \\
\midrule
simulation-based search & busca baseada em simulação
  & branching factor      & fator de ramificação \\
decision-time planning  & planejamento em tempo de decisão
  & anytime algorithm     & algoritmo \emph{anytime} \\
forward search          & busca dirigida para frente
  & best-first search     & busca best-first \\
expectimax tree         & árvore expectimax
  & progressive widening  & alargamento progressivo \\
rollout                 & \emph{rollout}, simulação até o fim
  & self-play             & auto-jogo \\
tree policy             & política de árvore
  & policy prior          & \emph{prior} de política \\
selection / expansion   & seleção / expansão
  & policy improvement    & melhoria de política \\
backup                  & retropropagação
  & expert iteration      & iteração por especialista \\
visit count             & contagem de visitas
  & latent state          & estado latente \\
upper confidence bound  & limite superior de confiança
  & test-time compute     & computação em tempo de decisão \\
\bottomrule
\end{tabular}}

\begin{alertabox}
\textbf{Erros comuns.} Devolver $\argmax\Qh$ em vez de $\argmax N$; esquecer de
inverter o sinal do retorno em jogos de dois jogadores; usar $c$ sem normalizar
os retornos; descartar a árvore a cada jogada; treinar $\bm{p}$ contra o
$\argmax$ da busca em vez da distribuição de visitas; supor que um modelo com
$1\%$ de erro por passo serve para um horizonte de $50$ passos.
\end{alertabox}

\section{Para ir além}

{\small
\begin{itemize}
  \item L. Kocsis e C. Szepesvári. Bandit based Monte-Carlo planning.
        \emph{ECML}, 2006. --- o artigo original do UCT.
  \item C. Browne et al. A survey of Monte Carlo tree search methods.
        \emph{IEEE Transactions on CIAIG}, 4(1):1--43, 2012.
  \item P.-A. Coquelin e R. Munos. Bandit algorithms for tree search.
        \emph{UAI}, 2007. --- os limites de pior caso do UCT.
  \item D. Silver et al. Mastering the game of Go with deep neural networks and
        tree search. \emph{Nature}, 529:484--489, 2016.
  \item D. Silver et al. Mastering the game of Go without human knowledge.
        \emph{Nature}, 550:354--359, 2017.
  \item D. Silver et al. A general reinforcement learning algorithm that masters
        chess, shogi, and Go through self-play. \emph{Science},
        362:1140--1144, 2018.
  \item T. Anthony, Z. Tian e D. Barber. Thinking fast and slow with deep
        learning and tree search. \emph{NeurIPS}, 2017.
  \item J. Schrittwieser et al. Mastering Atari, Go, chess and shogi by planning
        with a learned model. \emph{Nature}, 588:604--609, 2020.
  \item R. Sutton e A. Barto. \emph{Reinforcement Learning: An Introduction},
        2ª ed., MIT Press, 2018. Capítulo 8.
  \item E. Brunskill. \emph{CS234: Reinforcement Learning}, Stanford University
        --- \emph{Lecture 13: Monte Carlo Tree Search}, com material derivado de
        David Silver.
\end{itemize}}

\vfill
{\footnotesize\color{rlCinza}
Material adaptado e expandido de \textbf{CS234: Reinforcement Learning},
Emma Brunskill, Stanford University. Prof.\ Marcelo Keese Albertini,
Faculdade de Computação, Universidade Federal de Uberlândia.}

\end{document}
